% Generated draft: requires mathematical and source review.
\documentclass[11pt]{article}
\usepackage[margin=25mm]{geometry}
\usepackage{fontspec}
\setmainfont{Latin Modern Roman}
\usepackage{amsmath,amssymb,mathtools,unicode-math}
\setmathfont{Latin Modern Math}
\usepackage{enumitem,needspace,xurl,hyperref}
\setlength{\parindent}{0pt}
\setlength{\parskip}{5pt}
\setlength{\emergencystretch}{3em}
\newcommand{\N}{\mathbb{N}}
\newcommand{\Z}{\mathbb{Z}}
\newcommand{\Q}{\mathbb{Q}}
\newcommand{\R}{\mathbb{R}}
\newcommand{\C}{\mathbb{C}}
\newcommand{\F}{\mathbb{F}}
\newcommand{\A}{\mathbb{A}}
\newcommand{\PP}{\mathbb{P}}
\newcommand{\E}{\mathbb{E}}
\newcommand{\eps}{\varepsilon}
\newcommand{\dd}{\,\mathrm{d}}
\begin{document}
\section*{UnsolvedMath ID: 800024}
\textbf{Source problem code:} catalogue entry
\subsection{Definitions and mathematical statement}
A semigroup is a nonempty set equipped with an associative binary operation. For a mathematical structure \(A\), let \(\operatorname{End}(A)\) denote its semigroup of endomorphisms under composition. The meaning of endomorphism depends on the specified type of structure. For a partially ordered set \((A,\leq_A)\), the usual convention is that an endomorphism is a map \(f:A\to A\) satisfying \(x\leq_A y\Rightarrow f(x)\leq_A f(y)\). For an algebra in a fixed signature, endomorphisms preserve all its specified operations. A relatively free algebra on a set \(X\) in a variety \(\mathcal V\) is an algebra \(F\in\mathcal V\) with a generating map \(i:X\to F\) such that every map \(g:X\to B\), for \(B\in\mathcal V\), extends uniquely to a homomorphism \(\bar g:F\to B\) satisfying \(\bar g\circ i=g\).

To express the reconstruction question precisely once its parameters have been specified, let \(\mathcal D\) be a comparison class of structures of a specified type, let \(\mathcal K\subseteq\mathcal D\), and associate to each \(A\in\mathcal D\) a specified nonempty subsemigroup \(S_A\leq\operatorname{End}(A)\). Say that this association determines the members of \(\mathcal K\) up to isomorphism relative to \(\mathcal D\) if
\[
\forall A\in\mathcal K\;\forall B\in\mathcal D,\qquad
S_A\cong_{\mathrm{sgp}}S_B\ \Longrightarrow\ A\cong_{\mathrm{str}}B.
\]
Here \(\cong_{\mathrm{sgp}}\) means abstract semigroup isomorphism, and \(\cong_{\mathrm{str}}\) means isomorphism of the underlying structures in the specified type. The special case \(S_A=\operatorname{End}(A)\) concerns reconstruction from the full endomorphism semigroup.

\textbf{Problem attributed to E. S. Lyapin (Problem 25; permanent catalogue entry).} For natural classes of mathematical structures, including ordered sets and relatively free algebras, determine the subclasses whose members are determined up to isomorphism by their endomorphism semigroups or specified subsemigroups thereof, and determine which such subsemigroups have this reconstruction property.

This is a research programme rather than a single fully specified conjecture. The displayed implication supplies a parameterized formulation; the source does not fix the comparison class \(\mathcal D\), the target subclass \(\mathcal K\), the algebraic signatures or varieties, or the choices \(S_A\). In particular, reconstruction within a subclass corresponds to taking \(\mathcal D=\mathcal K\), whereas reconstruction against a larger ambient class is a stronger requirement. The source does not distinguish these interpretations or specify whether abstract semigroup isomorphisms alone are intended. These choices remain unresolved.

The supplied background describes connections between endomorphism semigroups and underlying structures as a direction emphasized by Lyapin in his book \emph{Semigroups}, and as a natural generalization of ideas of Galois theory. It also reports that Lyapin obtained reconstruction results for various classes of partially ordered sets, citing [Lyapin4], [L5], and [L6]; their precise hypotheses and conclusions are not supplied here.
\subsection{Short English statement}
Problem 25, attributed to E. S. Lyapin, asks which natural classes of structures can be reconstructed up to isomorphism from their endomorphism semigroups or selected subsemigroups. Examples include ordered sets and relatively free algebras. The comparison classes and permitted subsemigroups must be specified before this becomes a definite reconstruction problem. The registry labels the problem open; this is a supplied status, not an independently verified assessment.
\subsection{Sources}
\begin{itemize}
\item Local UnsolvedMath registry, supplied complete record: permanent ID \(800024\), source code \texttt{catalogue entry}, title \emph{The Lyapin's Notebook --- Problem 25}, AMR Open Problem Lists. Upstream revision: \texttt{d371ee917ac9647fbb06cd6fc529deec6f3dda82}. The supplied record labels the problem \emph{open}.
\item Bershadsky, Kublanovsky, and Mashevitzky, \emph{The Lyapin's Notebook: Unsolved Problems in Semigroup Theory} (2026), Problem 25. Supplied source URL: \url{https://arxiv.org/abs/2604.04763}. The record reports access on 2026-10-16 and extraction from \texttt{arxiv\_latex}; these are supplied provenance details, not an independently verified source visit.
\item The statement attributes the problem to E. S. Lyapin. The registry's proposer fields list Bershadsky, Kublanovsky, and Mashevitzky and the year 2026; these collection-level fields do not establish the date when Lyapin originally proposed the problem.
\item E. S. Lyapin, \emph{Semigroups}, mentioned in the supplied statement without edition, publication year, or page reference. The supplied statement also cites [Lyapin4], [L5], and [L6] for reconstruction results concerning partially ordered sets, but provides no corresponding bibliographic entries.
\end{itemize}
\end{document}
